\conclusion

В результате выполнения выпускной квалификационной работы была спроектирована и реализована backend-часть открытой микросервисной платформы CoachLink, предназначенной для автоматизации взаимодействия в системе <<тренер-спортсмен>>. Разработанная система решает практическую проблему разрозненного обмена тренировочными планами и отчётами через мессенджеры, электронные таблицы и произвольные файлы, предоставляя единую серверную платформу для структурированного хранения и обработки тренировочных данных.

Целью работы была разработка открытой микросервисной backend-платформы для стандартизации и автоматизации взаимодействия между тренерами и спортсменами. Данная цель была достигнута за счёт проектирования и реализации набора независимых сервисов, каждый из которых отвечает за отдельную часть логики предметной области. Реализованная платформа охватывает полный цикл базового взаимодействия: от регистрации пользователя и установления связи между тренером и спортсменом до назначения тренировки, отправки отчёта, уведомления участников и последующего анализа данных.

В ходе выполнения работы были последовательно решены все поставленные задачи. В аналитической главе проведён обзор существующих способов организации тренировочного процесса. Показано, что мессенджеры, электронные таблицы, спортивные трекеры и закрытые платформы решают отдельные части проблемы, но не предоставляют одновременно единый формат отчётности, поддержку ролей, групповые назначения, открытую архитектуру и возможность самостоятельного развёртывания. На основе анализа сформулированы функциональные и нефункциональные требования к backend-платформе.

В проектной части предложена микросервисная архитектура, в которой функциональность разделена между API Gateway, Auth Service, User Service, Training Service, Notification Service, Analytics Service и AI Service. Для межсервисного взаимодействия используются синхронные HTTP-вызовы и событийная коммуникация через NATS JetStream. Спроектирована модель данных с разделением баз по сервисам, определены ключевые пользовательские сценарии и описан публичный REST API в формате OpenAPI.

В главе реализации описана практическая разработка backend-сервисов на языке Go. Каждый сервис имеет собственную область ответственности и собственную схему данных. Stateless-сервисы, такие как Analytics Service и AI Service, работают как вычислительные модули поверх данных других сервисов. Проведено тестирование на трёх уровнях: unit-тесты бизнес-логики, E2E smoke-сценарий и интеграционные тесты через API Gateway. Все тестовые проверки выполнены успешно, что подтверждает работоспособность как отдельных сервисов, так и межсервисного взаимодействия.

Практическая значимость работы заключается в том, что разработанный backend предоставляет тренеру единый механизм назначения тренировок и отслеживания их выполнения, а спортсмену --- стандартизированный формат отчёта, пригодный для последующего анализа. Платформа формирует централизованную историю тренировочного процесса, поддерживает систему уведомлений о ключевых событиях и содержит вспомогательный AI-модуль для формирования рекомендаций по последним отчётам спортсмена. Благодаря открытому API система может быть использована различными клиентскими приложениями без привязки к конкретному пользовательскому интерфейсу.

Архитектурным результатом работы является разделение системы на независимые сервисы с чёткими зонами ответственности. Такое разделение повышает сопровождаемость: изменение механизма уведомлений не затрагивает сервис тренировок, развитие аналитики не влияет на авторизацию, а расширение AI-сценариев может выполняться отдельно от базового цикла назначений и отчётности.

Дальнейшее развитие платформы может включать поддержку других видов спорта за счёт расширения модели тренировочных планов и отчётов, интеграцию со спортивными устройствами для автоматического заполнения части отчёта, развитие аналитики с анализом динамики нагрузки и выявлением тенденций, а также расширение AI-сценариев. Предметная реализация первой версии ориентирована на лёгкую атлетику, однако архитектура платформы спроектирована с учётом этих направлений и не требует полного перепроектирования для их реализации.

Поставленная цель достигнута, а задачи выпускной квалификационной работы выполнены в полном объёме.
